\documentclass[11pt]{article}
\usepackage[a4paper,margin=28mm]{geometry}
\usepackage[T1]{fontenc}
\usepackage{lmodern,amsmath,amssymb,amsthm,mathtools,microtype}
\usepackage{xurl,hyperref,enumitem,booktabs}
\hypersetup{colorlinks=true,linkcolor=blue,citecolor=blue,urlcolor=blue,
 pdftitle={Pinned distance densities in nested train track models},
 pdfauthor={AI-assisted manuscript candidate}}
\setlength{\emergencystretch}{2em}
\numberwithin{equation}{section}
\newtheorem{theorem}{Theorem}[section]
\newtheorem{lemma}[theorem]{Lemma}
\newtheorem{proposition}[theorem]{Proposition}
\newtheorem{corollary}[theorem]{Corollary}
\theoremstyle{definition}\newtheorem{definition}[theorem]{Definition}
\theoremstyle{remark}\newtheorem{remark}[theorem]{Remark}
\newcommand{\R}{\mathbb R}
\newcommand{\N}{\mathbb N}
\newcommand{\supp}{\operatorname{supp}}
\newcommand{\Var}{\operatorname{Var}}
\newcommand{\dist}{\operatorname{dist}}
\newcommand{\Unif}{\operatorname{Unif}}
\newcommand{\Bad}{\operatorname{Bad}}
\newcommand{\dimH}{\dim_{\mathrm H}}
\newcommand{\dimP}{\dim_{\mathrm P}}
\newcommand{\dimB}{\overline{\dim}_{\mathrm B}}
\newcommand{\norm}[1]{\lVert #1\rVert}
\newcommand{\abs}[1]{\lvert #1\rvert}
\newcommand{\ind}{\mathbf 1}
\title{Pinned distance densities in nested train track models}
\author{AI-assisted manuscript candidate}
\date{8 October 2026}
\begin{document}
\maketitle
\begin{abstract}
We study the original pinned distance probabilities of explicit deterministic
planar product-digit measures with singular coordinate marginals. For every
rational $1<\alpha<2$, a superlacunary construction has Hausdorff dimension
$\alpha$ and packing dimension two. Its distance densities have the uniform
attained exponent $q_*=(2-\alpha)/(3-2\alpha)$ when $\alpha<3/2$, and are
bounded and jointly continuous when $\alpha\geq3/2$. The finite exponent is
sharp precisely at pins whose first coordinate belongs to the horizontal
support; off that support the densities are locally bounded and continuous.
At a fixed stage ratio the critical exponent changes and its finite endpoint
fails. We determine the corresponding exceptional pin sets and their
Hausdorff and packing dimensions. A two-step bounded-variation and coarea
estimate gives a fixed positive restriction with a jointly continuous density
at every pin, and hence intervals of uniformly positive length in all the
studied distance sets. At the fixed-ratio transition the densities belong to
every finite $L^q$, with an $O(\log q)$ norm upper bound and a
double-exponential tail upper bound, but have no continuous version on the
horizontal-support pin set. Essential boundedness there remains open.
These conclusions use the exact rectangular placement and stage schedule;
they are not general Falconer-type dimension thresholds.
\end{abstract}
\setcounter{tocdepth}{1}
\tableofcontents
\section{Introduction and main results}\label{sec:introduction}
For a pin $p\in\R^2$, write $D_p(y)=|y-p|$. The size of the compact
set $D_p(E)$ and the regularity of a particular probability
$(D_p)_*\mu$ supported on it are different questions. This paper gives
explicit nested examples in which every studied distance set contains an
interval, while the original probability can fail $L^2$ at every original
pin. It also shows how changing the stage schedule changes both the critical
integrability exponent and whether that exponent is attained.

Throughout, $\alpha\in(1,2)$ is rational, and
\begin{equation}\label{eq:parameters}
 a=\frac{\alpha-1}{2},\qquad b=\frac\alpha2,\qquad
 c=1-b,\qquad \beta=\alpha-1.
\end{equation}
In particular $b-a=1/2$ and $a+c=1/2$. Section~\ref{sec:construction}
constructs compact sets $E_-$ and $E_+$ and probabilities
$\lambda=\kappa\times\zeta_-$ and $\mu=\kappa\times\zeta_+$ on them.
For a fixed $\eta=2^{-k}<1/8$, they lie respectively in
\[
 P_-=[0,\eta]^2,\qquad P_+=[0,\eta]\times[1-\eta,1].
\]
Pins range over the whole rectangle $P=P_-$. Set
\[
 H=\supp\kappa,\qquad \mathcal B=H\times[0,\eta],\qquad
 \nu_p=(D_p)_*\mu.
\]
Thus $E_-\subset\mathcal B\subset P$. When $\nu_p$ is absolutely
continuous, $g_p$ denotes its density. All $L^q$ norms below use Lebesgue
measure on $\R$. A continuous density means a continuous version on all of
$\R$, including the endpoints of its support.

\begin{theorem}[Superlacunary schedule]\label{thm:main-super}
Use the digit construction of Section~\ref{sec:construction} with
$L_j/n_{j-1}\to\infty$. Then both $E_\pm$ have Hausdorff dimension
$\alpha$ and packing dimension two, and both $\lambda,\mu$ are
$\alpha$-Frostman probabilities with singular coordinate marginals.
Every $\nu_p$, $p\in P$, is absolutely continuous.
\begin{enumerate}[label=(\roman*),leftmargin=*]
\item If $1<\alpha<3/2$, put
\begin{equation}\label{eq:qstar}
 q_* = \frac{2-\alpha}{3-2\alpha}=\frac{c}{c-a}.
\end{equation}
Then $\sup_{p\in P}\norm{g_p}_{q_*}<\infty$. For every
$p\in\mathcal B$, the exact membership range is $1\leq q\leq q_*$;
no such density belongs to any $L^q$ with $q>q_*$, including $q=\infty$.
\item If $3/2\leq\alpha<2$, the densities admit one bounded jointly
continuous version on $P\times\R$.
\item In every parameter range, the original density is locally jointly
continuous on $(P\setminus\mathcal B)\times\R$, and
$\norm{g_p}_\infty\leq C/\dist(p_1,H)$ for $p\notin\mathcal B$.
\end{enumerate}
\end{theorem}
At $\alpha=4/3$, the uniform upper conclusion is $L^2$ for \emph{all}
$p\in P$; failure of every $L^{2+\varepsilon}$ applies exactly to pins
with $p_1\in H$. There is no conflict with boundedness off $H$.
The endpoint $\alpha=3/2$ is a direct uniform-convergence theorem,
not a formal limiting value of the finite exponent.

\begin{theorem}[Fixed stage ratio]\label{thm:main-fixed}
Use $L_j=K n_{j-1}$ for $j\geq2$, with an integer $K\geq1$, and set
\begin{equation}\label{eq:DK-main}
 D_K=K(3/2-\alpha)+1-\alpha,
 \qquad \beta_K=\frac{\alpha+2aK}{1+aK}.
\end{equation}
The sets $E_\pm$ have Hausdorff dimension $\alpha$ and packing dimension
$\beta_K<2$, and their probabilities are $\alpha$-Frostman with singular
coordinate marginals. Assume $aK>2$ for the following regularity statements.
\begin{enumerate}[label=(\roman*),leftmargin=*]
\item If $D_K>0$, set $q_K=c(K+1)/D_K$. For every $1\leq q<q_K$,
$\sup_{p\in P}\norm{g_p}_q<\infty$. At every $p\in\mathcal B$, the
exact range is $1\leq q<q_K$: the finite critical endpoint fails.
\item If $D_K<0$, the original densities are bounded and jointly
continuous on $P\times\R$.
\item If $D_K=0$, all finite $L^q$ upper bounds hold, and
\begin{equation}\label{eq:boundary-main}
 \sup_{p\in P}\norm{g_p}_q\leq C\log(e+q),\qquad 1\leq q<\infty.
\end{equation}
For each $p\in\mathcal B$ there is no continuous-density version.
Whether $g_p\in L^\infty$ at such pins remains open. Equation
\eqref{eq:boundary-main} is only an upper bound; no matching norm growth
lower bound is asserted.
\item Off $\mathcal B$, the local bounded-continuity conclusion of
Theorem~\ref{thm:main-super}(iii) holds in all three regimes.
\end{enumerate}
\end{theorem}

\begin{theorem}[Positive restrictions and actual distance intervals]
\label{thm:main-intervals}
Assume either the superlacunary schedule, or the fixed-ratio schedule with
$aK>1$. For each fixed $0<\varepsilon\leq1$, the positive subprobability
\[
 \nu_{p,\varepsilon}=(D_p)_*\bigl(
 \ind_{\{|y_1-p_1|\geq\varepsilon\}}\mu\bigr)
\]
has a jointly continuous density $g_\varepsilon(p,r)$ satisfying
\[
 0\leq g_\varepsilon(p,r)\leq C/\varepsilon,
 \qquad \nu_{p,\varepsilon}(\R)\geq1-C\varepsilon^{\alpha-1}.
\]
There exists $\ell>0$ such that every $D_p(E_+)$, $p\in P_-$, contains
an interval of length $\ell$. The reversed assertion holds from $P_+$
to $E_-$. In particular every pin in $E_-\cup E_+$ has an interval in
its distance set to that union. The interval may depend on the pin;
only its length lower bound is uniform.
\end{theorem}

We also determine the exact exceptional sets for stronger density
regularity, their dimensions, and positive-mask stability. At the
fixed-ratio boundary we obtain a double-exponential \emph{upper} tail and
uniform double-exponential integrability on the bounded output interval.

\subsection{Relation to earlier work and verification status}
The near-track obstruction and far-track smoothing distinction already
appears in Guth--Iosevich--Ou--Wang \cite[Sections~1.2 and~6]{GIOW}.
Their general good-source construction uses a wave-packet decomposition.
Here the restriction is an actual positive, fixed horizontal-band deletion
for one explicitly placed nested model. Neither that finite-scale intuition
nor a general distance-set threshold is claimed as new. Peres--Schlag
\cite{PS} provides classical generalized-projection background; no result
from that theory is an input to the proofs below.

This candidate consolidates four analytic source notes after a separate
model-based analytic review and a limited recheck of two scope
clarifications. That process is not professional human peer review or
formal verification. The consolidation is expository, and no priority claim
or comprehensive literature-novelty certification is made. The package's
source map records the exact inputs and editorial changes. All mathematical
arguments needed here are included in the paper.

\section{The digit construction and dimensions}\label{sec:construction}
Choose positive integers $L_j$ with $aL_j$, $L_j/2$, and $bL_j$ integral.
Let $n_0=0$ and $n_j=\sum_{i=1}^j L_i$. At stage $j$, the free horizontal
binary positions, counted relative to $n_{j-1}$, are
\begin{equation}\label{eq:free-positions}
 \{1,\ldots,aL_j\}\ \cup\ \{L_j/2+1,\ldots,L_j\};
\end{equation}
the free vertical positions are $\{1,\ldots,bL_j\}$. All remaining
digits are zero. All free bits in both coordinates are independent and fair.
Let $\kappa_X$ and $\kappa_Y$ be the resulting probabilities.

Fix $k$ with $\eta=2^{-k}<1/8$ and take $L_1$ sufficiently large that
the first $k$ positions are free in both coordinates. Condition $\kappa_X$
on its all-zero length-$k$ prefix to obtain $\kappa$. Condition $\kappa_Y$
on its all-zero and all-one length-$k$ prefixes to obtain $\zeta_-$ and
$\zeta_+$ respectively. Define
\begin{equation}\label{eq:measures}
 \lambda=\kappa\times\zeta_-,\quad
 \mu=\kappa\times\zeta_+,\quad E_\pm=\supp(\kappa\times\zeta_\pm).
\end{equation}
Supports are images of the compact coding spaces. At a binary endpoint one
uses a supporting closed coding cylinder; no zero-mass endpoint is added
by an arbitrary half-open convention. The fixed prefix conditioning changes
mass constants by $2^k$ in each coordinate, and never changes a later
stage exponent.

For every $p\in P_-$ and $y\in P_+$, including filled-tail points, set
\begin{equation}\label{eq:separation}
 v_0=1-2\eta>0,\qquad D_0=2;
 \qquad y_2-p_2\geq v_0,\quad |y-p|\leq D_0.
\end{equation}
Hence the vertical derivative of distance is at least $v_0/D_0$.
All output probabilities are supported in $I=[v_0,D_0]$.

\subsection{Two exact schedules}
In the superlacunary schedule, assume
\begin{equation}\label{eq:superlacunary}
 \frac{L_j}{n_{j-1}}\longrightarrow\infty\qquad(j\geq2).
\end{equation}
For any rational $\alpha$, choose an integer $d$ for which $ad,bd,d/2$
are integers, choose an admissible large $L_1$, and take
$L_j=d(n_{j-1}+1)^2$ for $j\geq2$. This gives an actual realization.
For example $\alpha=13/10$, $\eta=1/16$, $L_1=40$, and
$L_j=20(n_{j-1}+1)^2$ give $a=3/20$, $b=13/20$ and
$n_1-aL_2=-5003$.

In the fixed-ratio schedule, an integer $K\geq1$ is fixed and
\begin{equation}\label{eq:fixed-schedule}
 L_j=K n_{j-1},\quad n_j=(K+1)n_{j-1}\quad(j\geq2).
\end{equation}
Choosing $L_1$ divisible by the necessary denominators preserves all
integrality requirements. In particular $n_{j-1}=n_1(K+1)^{j-2}$.
No rounding argument or assertion for irrational $\alpha$ is used.

The original non-$L^2$ example was restricted to $1<\alpha<4/3$ and also
imposed $(2-3\alpha/2)L_j-\alpha n_{j-1}\to\infty$.
Condition~\eqref{eq:superlacunary} implies that condition in this range.
There is no need to impose it for the enlarged family $1<\alpha<2$.

\subsection{Frostman bounds and Hausdorff dimension}
Let $F_X(N),F_Y(N)$ count free positions through depth $N$, before prefix
conditioning. In a stage of length $L$ the relative counts are
\[
 f_X(l)=\min(l,aL)+\max(0,l-L/2),\qquad f_Y(l)=\min(l,bL).
\]
Their sum equals
\begin{equation}\label{eq:total-free}
 f_X(l)+f_Y(l)=
 \begin{cases}
 2l,&0\leq l\leq aL,\\
 l+aL,&aL\leq l\leq L/2,\\
 2l-cL,&L/2\leq l\leq bL,\\
 l+(\alpha-1)L,&bL\leq l\leq L.
 \end{cases}
\end{equation}
Each expression is at least $\alpha l$, and the endpoint value is
$\alpha L$. Consequently
\begin{equation}\label{eq:counts}
 F_X(N)+F_Y(N)\geq\alpha N,\qquad
 F_X(n_j)=F_Y(n_j)=bn_j.
\end{equation}
For $N\geq k$, every occupied depth-$N$ square has mass
$2^{2k-F_X(N)-F_Y(N)}$. A ball of radius comparable to $2^{-N}$ meets
only a bounded number of such squares. Thus both probabilities obey the
$\alpha$-Frostman estimate $\rho(B(x,r))\leq C r^\alpha$.
At $N=n_j$ there are at most $C2^{\alpha n_j}$ occupied squares, so
$\dimH E_\pm\leq\alpha$. The Frostman estimate gives the reverse
inequality and hence equality.

Also $f_X(l)\geq\beta l$ and $b\geq\beta$, whence
\begin{equation}\label{eq:horizontal-frostman}
 F_X(N)\geq\beta N,\qquad
 \kappa([u-\varepsilon,u+\varepsilon])\leq K_X\varepsilon^\beta
 \quad(0<\varepsilon\leq1).
\end{equation}
The interval bound follows by comparison with dyadic intervals and includes
all real $u$. In particular $\kappa$ is nonatomic. Each coordinate support
has a cover at depth $n_j$ by $C2^{bn_j}$ intervals of length $2^{-n_j}$.
The total length tends to zero because $b<1$. The coordinate laws are
therefore singular, and are nonatomic because their free counts diverge.
There is no hidden Lebesgue coordinate in the limiting product.

\subsection{Packing dimensions of the planar supports}
Under~\eqref{eq:superlacunary}, in each sufficiently late stage the first
$aL_j$ bits of both coordinates are free. Every fixed coding cylinder thus
contains at least $2^{2aL_j}$ occupied squares at depth
$n_{j-1}+aL_j$, and
\[
 \frac{2aL_j}{n_{j-1}+aL_j}\longrightarrow2.
\]
Every such cylinder has upper box dimension two. Every nonempty relatively
open subset of the support contains a coding cylinder. To pass from upper
box to packing dimension, use the countable-cover characterization by
modified upper box dimension: take relative closures in any countable
cover, which do not change upper box dimension. Baire's theorem implies
that one closure contains a relatively open subset. Its upper box dimension
is two, proving $\dimP E_\pm=2$.

For~\eqref{eq:fixed-schedule}, the total-count ratio during a stage is
$[\alpha n+f_X(l)+f_Y(l)]/(n+l)$, with $L=Kn$. It is monotone on each
of the four pieces in~\eqref{eq:total-free}. Its values at $l=0,L/2,L$
are $\alpha$, and its two possible maxima are
\[
 B_1=\frac{\alpha+2aK}{1+aK},\qquad
 B_2=\frac{\alpha+(3b-1)K}{1+bK}.
\]
After clearing denominators, the numerator of $B_1-B_2$ is
$a(1-b)K^2>0$. Thus the upper box dimension is
$\beta_K=B_1<2$. A fixed coding prefix changes the free count by a constant,
so every cylinder has this same upper box dimension. The preceding Baire
argument proves
\begin{equation}\label{eq:planar-dimensions}
 \dimH E_\pm=\alpha,\qquad
 \dimP E_\pm=
 \begin{cases}2,&\text{superlacunary},\\
 (\alpha+2aK)/(1+aK),&\text{fixed ratio}.
 \end{cases}
\end{equation}
No regular subset is substituted for the original supports.

\subsection{Filled tails and the exact replacement}
Choose a stage $j_0$ with $n_{j_0}\geq k$. For $j\geq j_0$, define
$\mu_j$ by retaining the first $n_j$ digits of $\mu$ and replacing all
later digits by independent fair bits. Equivalently, each occupied
$2^{-n_j}$-square is filled uniformly with its exact $\mu$-mass.
Then $\mu_j\rightharpoonup\mu$.

Inside a parent square $Q$ of side $h=2^{-n}$, where $n=n_{j-1}$,
the conditional parent law is uniform and the next law is the affine image
of $\xi_{A,w}\times\zeta_{B,v}$, where
\begin{align}
 \xi_{A,w}&=\frac1A\sum_{i=0}^{A-1}\Unif[i/A,i/A+w],\nonumber\\
 \zeta_{B,v}&=\frac1B\sum_{l=0}^{B-1}\Unif[l/B,l/B+v],\label{eq:template}\\
 A&=2^{aL},\quad B=2^{bL},\quad w=2^{-L/2},\quad v=2^{-L},
 \qquad L=L_j.\nonumber
\end{align}
The initial horizontal free block chooses the $A$ positions; the last half
of the horizontal bits and the filled tail produce the entire width $w$.
The vertical free block chooses $B$ positions and the filled tail supplies
width $v$. Therefore
\begin{equation}\label{eq:resolution}
 Bw=A,\qquad Bv=2^{-cL},\qquad w\leq A^{-1},\quad v\leq B^{-1}.
\end{equation}
This exact placement, rather than the number of descendants alone, is the
input to the next section.

\section{Two bounded variation quadrature estimates}\label{sec:quadrature}
We first isolate the deterministic geometric estimate used at every stage.
Let $U$ denote uniform probability on $[0,1]$.

\begin{lemma}[Equal-bin quadrature]\label{lem:BV}
If a probability $\xi$ assigns mass $1/N$ to each of the $N$ equal bins
of $[0,1]$, then for every function $F$ of bounded variation,
\begin{equation}\label{eq:BV}
 \left|\int F\,d\xi-\int F\,dU\right|\leq\frac{\Var(F)}N.
\end{equation}
\end{lemma}
\begin{proof}
Within each bin the difference between the two conditional averages is
bounded by the oscillation of $F$. Summing with weight $1/N$ proves the
claim, since the sum of the bin oscillations is at most total variation.
At jump endpoints use one-sided representatives consistent with the bin
convention. The template laws~\eqref{eq:template} used below are absolutely
continuous within bins, so endpoint values do not affect their integrals.
\end{proof}

Let $Q=[q_1,q_1+h]\times[q_2,q_2+h]$, $0<h\leq1$, satisfy
\eqref{eq:separation} with respect to the pin $p$. Let $U_Q$ be uniform
probability on $Q$, and let $\rho_Q$ be the affine image of the product
in~\eqref{eq:template}, allowing any integers $A,B\geq1$ and widths
$0<w\leq A^{-1}$, $0<v\leq B^{-1}$. Write
\[
 M_\varepsilon(p,y)=\ind_{\{|y_1-p_1|\geq\varepsilon\}}.
\]
When a norm of a signed pushforward is written below, it means the norm of
its density. Each constituent measure is positive and unnormalized.

\begin{lemma}[Fixed-band comparison]\label{lem:fixed-band}
Both masked distance pushforwards have densities, and
\begin{equation}\label{eq:fixed-band}
 \norm{(D_p)_*(M_\varepsilon\rho_Q)
       -(D_p)_*(M_\varepsilon U_Q)}_\infty
 \leq\frac{4D_0}{h}
       \left(\frac1{v_0A}+\frac1{\varepsilon Bw}\right).
\end{equation}
The pin and the mask are identical in the two terms.
\end{lemma}
\begin{proof}
Use coordinates $X=q_1+hu$, $Y=q_2+hz$ and the local output
$t=|y-p|/h$. Conversion back to physical radius contributes a final
factor $h^{-1}$ to the density.

First fix a horizontal strip $i/A\leq u\leq i/A+w$, keep the part
$|X-p_1|\geq\varepsilon$, and average in the vertical coordinate.
For a fixed local output $t$, set $r=ht$. Horizontal coarea gives the
$t$-density
\begin{equation}\label{eq:horizontal-coarea}
 G_i(z,t)=\frac1w\sum_{\text{allowed roots }X}\frac{r}{|X-p_1|},
 \qquad X-p_1=\pm\sqrt{r^2-(Y-p_2)^2}.
\end{equation}
Indeed $\partial t/\partial u=(X-p_1)/r$, so no extra factor of $h$
belongs in~\eqref{eq:horizontal-coarea}. On each branch, $Y-p_2>0$
and $|X-p_1|$ decreases as $z$ increases. The reciprocal derivative
$r/|X-p_1|$ is monotone. Requiring the root to lie in the strip cuts out
an interval of $z$, and imposing the band cutoff only shortens it.
The zero extension of each monotone interval function has variation at
most twice its supremum. There are at most two branches, each at most
$D_0/(\varepsilon w)$, so
\begin{equation}\label{eq:vertical-variation}
 \Var_z G_i(z,t)\leq\frac{4D_0}{\varepsilon w}.
\end{equation}
Lemma~\ref{lem:BV} replaces $\zeta_{B,v}$ by $U$ at a density cost
$4D_0/(\varepsilon Bw)$. Averaging the strips with their $1/A$ weights
preserves that bound. The width $v$ has not entered except through
$v\leq1/B$.

With the vertical coordinate uniform, fix $u$. There is only one allowed
vertical root because $Y-p_2\geq v_0$. Its local-output density is
\begin{equation}\label{eq:vertical-coarea}
 J(u,t)=\ind_{\{\text{allowed vertical root}\}}
        \ind_{\{|X-p_1|\geq\varepsilon\}}
        \frac{r}{Y-p_2}.
\end{equation}
It is at most $D_0/v_0$. On either side of $X=p_1$ it is monotone in
$|X-p_1|$; requiring the root to lie in the vertical side of $Q$ leaves
at most one interval per side. Therefore
$\Var_u J(u,t)\leq4D_0/v_0$. A second application of
Lemma~\ref{lem:BV}, this time to $\xi_{A,w}$, costs $4D_0/(v_0A)$.
Adding the two errors and returning to physical radius proves the claim.
Coarea is applied on strictly monotone intervals; the cutoff removes the
only horizontal critical point. All bounds hold for almost every output,
which is sufficient for the stated essential supremum.
\end{proof}

\begin{lemma}[Harmonic improvement]\label{lem:harmonic}
At the particular physical cutoff $\varepsilon=h/A$,
\begin{equation}\label{eq:harmonic}
 \norm{(D_p)_*(M_{h/A}\rho_Q)-(D_p)_*(M_{h/A}U_Q)}_\infty
 \leq C D_0\left(\frac1{hv_0A}
              +\frac{\log(2A)}{h^2Bw}\right).
\end{equation}
\end{lemma}
\begin{proof}
For the physical strip
$I_i=[q_1+hi/A,q_1+hi/A+hw]$, let $d_i$ be the infimum of
$|X-p_1|$ over its nonempty retained part. Empty strips are omitted.
Then $d_i\geq h/A$, and the sharper variation bound in the first part
of the preceding proof is $4D_0/(wd_i)$. Each strip lies in its own
coarse bin of length $h/A$. Order bins to the left and right of $p_1$.
Except for at most a fixed number of nearest neighbors, the $l$-th bin
has distance at least $lh/(2A)$. The cutoff controls the nearest bins.
The same ordering works if $p_1$ lies outside $Q$. Hence
\begin{equation}\label{eq:harmonic-sum}
 \frac1A\sum_i d_i^{-1}
 \leq\frac Ch\left(1+\sum_{l=1}^A\frac1l\right)
 \leq\frac Ch\log(2A).
\end{equation}
Average this sharper variation over strips before the vertical quadrature
step. Its local-output error is at most $CD_0\log(2A)/(hBw)$;
physical scaling gives the second term in~\eqref{eq:harmonic}.
The horizontal quadrature term is unchanged.
\end{proof}
The resolving relation $Bw=A\to\infty$ expresses a vertical sampling
spacing fine enough to resolve an actual horizontal interval. A Frostman
bound or a numerical branching profile does not imply either lemma.

\section{Positive filters and intervals at every pin}\label{sec:filters}
The arguments in this section apply whenever
\begin{equation}\label{eq:S}
 S=\sum_{j>j_0}2^{n_{j-1}-aL_j}<\infty.
\end{equation}
This follows from superlacunarity, since eventually $aL_j\geq2n_{j-1}$
and $n_{j-1}\geq j-1$. For a fixed ratio it follows from $aK>1$, since
the summands are $2^{-(aK-1)n_{j-1}}$.

\begin{proposition}[One fixed positive restriction]\label{prop:fixed-filter}
For every $p\in P$ and fixed $0<\varepsilon\leq1$, the actual measure
$\nu_{p,\varepsilon}=(D_p)_*(M_\varepsilon(p,\cdot)\mu)$ has a
jointly continuous density $g_\varepsilon(p,r)$, with
\begin{equation}\label{eq:fixed-good-bound}
 \norm{g_\varepsilon(p,\cdot)}_\infty\leq C_*/\varepsilon,
 \qquad \nu_{p,\varepsilon}(\R)\geq1-K_X\varepsilon^\beta.
\end{equation}
One possible constant is
\begin{equation}\label{eq:Cstar}
 C_*=(D_0/v_0)2^{n_{j_0}}+4D_0(1+v_0^{-1})S.
\end{equation}
It may depend on the finitely many initial stages, but is independent of
$p$ and any finest observation scale.
\end{proposition}
\begin{proof}
Let $g_{j,p,\varepsilon}$ be the masked density for $\mu_j$.
Apply Lemma~\ref{lem:fixed-band} in every parent square, multiply by
its exact parent mass, and add. The masses sum to one. By $Bw=A$,
\begin{equation}\label{eq:fixed-increment}
 \norm{g_{j,p,\varepsilon}-g_{j-1,p,\varepsilon}}_\infty
 \leq4D_0(v_0^{-1}+\varepsilon^{-1})2^{n_{j-1}-aL_j}.
\end{equation}
No count of cells appears. At the initial filled stage, vertical coarea
on each square gives
$\norm{g_{j_0,p,\varepsilon}}_\infty\leq(D_0/v_0)2^{n_{j_0}}$.
Thus the densities converge in $L^\infty$ uniformly in $p$, and the
bound~\eqref{eq:Cstar} follows.

For a fixed $p$, the discontinuities of the mask are on
$y_1=p_1\pm\varepsilon$. These lines have zero $\mu$-mass by
\eqref{eq:horizontal-frostman}. Weak convergence can therefore be applied
to a bounded continuous distance test function multiplied by the mask.
It identifies the weak limit with $(D_p)_*(M_\varepsilon\mu)$.
The common bounded output support makes the $L^\infty$ convergence also
$L^1$ convergence, identifying its density with the same measure.
This argument works for \emph{each} pin with the same bounds, without an
exceptional pin set. The retained-mass estimate follows directly from
\eqref{eq:horizontal-frostman}.

To obtain the continuous version, write the finite filled-tail density
$f_j$ as a finite linear combination of rectangle indicators. Polar coarea
gives, for $r>0$,
\begin{equation}\label{eq:polar-continuity}
 g_{j,\varepsilon}(p,r)=r\int_0^{2\pi}
 f_j(p+r\omega(\theta))
 \ind_{\{|r\cos\theta|\geq\varepsilon\}}\,d\theta,
 \quad\omega(\theta)=(\cos\theta,\sin\theta).
\end{equation}
A circle with positive radius meets each rectangle boundary in only finitely
many angles. The mask equality also has at most finitely many angles.
For a convergent sequence of $(p,r)$ the integrand therefore converges at
almost every angle and has a bounded integrable majorant. Dominated
convergence proves joint continuity. Around $r=0$ the density is zero by
uniform vertical separation; extend by zero for negative $r$.

For continuous radius fibers the supremum is the essential supremum.
Thus~\eqref{eq:fixed-increment} is a uniform supremum estimate on the
continuous versions over $P\times\R$. Their uniform limit is jointly
continuous and represents the actual positive restriction at every pin.
\end{proof}

\begin{corollary}[Raw absolute continuity and weak integrability]
\label{cor:weak}
Every $\nu_p$ is absolutely continuous. Its density satisfies
\begin{equation}\label{eq:weak-alpha}
 \sup_{p\in P}|\{r:g_p(r)>T\}|\leq C T^{-\alpha}
 \qquad(T>0).
\end{equation}
In particular $\sup_p\norm{g_p}_q<\infty$ for $1\leq q<\alpha$.
\end{corollary}
\begin{proof}
As $\varepsilon\downarrow0$, the positive restricted measures increase
to $\nu_p$, since $\mu\{y_1=p_1\}=0$. A Lebesgue-null set has zero
mass for every restriction and hence for the limit. This proves absolute
continuity. Write $g_p=g_\varepsilon(p,\cdot)+b_{p,\varepsilon}$ with
$b_{p,\varepsilon}\geq0$ and
$\int b_{p,\varepsilon}\leq K_X\varepsilon^\beta$.
For sufficiently large $T$, choose $\varepsilon=2C_*/T\leq1$.
Then the good density is at most $T/2$, so Markov's inequality gives
\[
 |\{g_p>T\}|\leq\frac2T\int b_{p,\varepsilon}
 \leq2K_X(2C_*)^\beta T^{-(1+\beta)}.
\]
Since $1+\beta=\alpha$, this is the claim for large $T$; bounded output
support handles smaller $T$. Layer-cake integration gives the stated
strong $L^q$ bounds.
\end{proof}

\begin{proposition}[Uniform length in the actual distance sets]
\label{prop:intervals}
There is $\ell>0$ such that every $D_p(E_+)$, $p\in P_-$, contains
an interval of length $\ell$. The reversed conclusion holds for
$p\in P_+$ and source $E_-$.
\end{proposition}
\begin{proof}
Choose a fixed $\varepsilon_0>0$ with
$K_X\varepsilon_0^\beta\leq1/2$. Every nonnegative fiber of
$g_{\varepsilon_0}$ then has mass at least $1/2$ on the interval
$I=[v_0,D_0]$ of length $W=D_0-v_0$. It has some point $r_p$ with
$g_{\varepsilon_0}(p,r_p)\geq1/(2W)$. Joint continuity and compactness,
using the continuous zero extension outside $I$, supply one $\delta>0$
such that
\[
 |r-r_p|\leq\delta\quad\Longrightarrow\quad
 g_{\varepsilon_0}(p,r)\geq1/(4W)
\]
for every $p$. The support of this submeasure is contained in the compact
set $D_p(E_+)$ because the submeasure is a positive restriction of the
actual source law. Thus the entire interval is in the actual image, not
only in a finite approximation or an almost-everywhere version of that
image. Taking any fixed $\ell\leq\delta$ proves the first assertion.

In the reverse configuration the vertical separation has negative sign,
but the absolute derivative bound is unchanged. The vertical branch is
still unique and the BV monotonicities simply reverse sign. The same proof
therefore applies. If necessary take the smaller of the two length bounds.
\end{proof}
Proposition~\ref{prop:intervals} proves Theorem~\ref{thm:main-intervals}.
In particular, failure of $L^2$ or even unboundedness of the raw density
is compatible with an interval in every one of the distance images here.

\section{The attained superlacunary endpoint}\label{sec:superlacunary}
We now use the exact global product structure to strengthen
Corollary~\ref{cor:weak}. Write $g_{j,p}$ for the unmasked filled-tail
distance density. In a stage beginning at $n=n_{j-1}$, write $L=L_j$,
$h=2^{-n}$ and use the parameters in~\eqref{eq:template}. Decompose the
\emph{signed increment} $g_{j,p}-g_{j-1,p}$ using the auxiliary band
\begin{equation}\label{eq:aux-band}
 \varepsilon_j=h/A.
\end{equation}
This changing band is only a proof device; it is not the fixed positive
restriction of Section~\ref{sec:filters}.

\subsection{The outer and inner parts of a stage increment}
Lemma~\ref{lem:harmonic} and $Bw=A$ give, after weighting by parent
masses, an outer-part bound
\begin{equation}\label{eq:outer}
 \norm{\text{outer increment}}_\infty
 \leq C(1+L)2^{2n-aL}.
\end{equation}
It also bounds every finite $L^q$ norm up to a constant depending only on
$I$. Under superlacunarity the series is summable: eventually
$\log_2(1+L)\leq aL/2$ and $aL/2\geq3n$, leaving $C2^{-n}$.

For the inner part it is essential to estimate the \emph{global} horizontal
mass. At depth $n$, every occupied horizontal prefix interval has
$\kappa$-mass $2^{k-bn}$. The next-stage coarse bins of width $h/A$
are globally aligned; each occupied strip has mass $2^{k-bn}/A$.
A band of length $2h/A$ meets only a bounded number of these bins.
For both the new filled law and the old filled law this gives
\begin{equation}\label{eq:global-M}
 M_j,M_{j-1}\leq C_k A^{-1}2^{-bn}.
\end{equation}
For the old law this is also immediate from its piecewise constant
horizontal density and the band's length. The corresponding vertical
densities satisfy
\begin{equation}\label{eq:global-H}
 H_j=2^{k+c(n+L)},\qquad H_{j-1}=2^{k+cn}.
\end{equation}
Horizontal restriction does not change the vertical conditional law:
both measures are actual product measures.

For a product submeasure of horizontal mass $M$ with vertical density
at most $H$, vertical coarea gives a distance-density supremum at most
$(D_0/v_0)MH$. Its total mass is $M$. Hence, putting
$\sigma=1-1/q$, with $1/\infty=0$,
\begin{equation}\label{eq:MH}
 \norm{\text{band density}}_q\leq C M H^\sigma,
 \qquad 1\leq q\leq\infty.
\end{equation}
For $q=\infty$ this is the direct supremum bound. Substituting
\eqref{eq:global-M}--\eqref{eq:global-H} gives for the new law
\begin{equation}\label{eq:inner-general}
 C2^{[-b+c\sigma]n}2^{[-a+c\sigma]L}.
\end{equation}
The old law has the additional favorable factor $2^{-c\sigma L}$.
The norm of their difference is bounded by the sum of these positive norms.

\begin{proposition}[Uniform endpoint upper bound]\label{prop:super-upper}
If $1<\alpha<3/2$, then
$\sup_{p\in P}\norm{g_p}_{q_*}<\infty$ for $q_*$ in~\eqref{eq:qstar}.
If $3/2\leq\alpha<2$, then the original densities are uniformly bounded
and have a jointly continuous version on $P\times\R$.
\end{proposition}
\begin{proof}
For $\alpha<3/2$, the endpoint identity is
\[
 c(1-1/q_*)=a.
\]
Thus~\eqref{eq:inner-general} reduces exactly to
$C2^{-(b-a)n}=C2^{-n/2}$. For $\alpha\geq3/2$, take $q=\infty$,
so $\sigma=1$ and the same bound becomes
\begin{equation}\label{eq:infinity-endpoint}
 C2^{-(\alpha-1)n}2^{-(\alpha-3/2)L}\leq C2^{-n/2}.
\end{equation}
In either case the inner parts are summable. Together with
\eqref{eq:outer}, this shows
\begin{equation}\label{eq:super-increment}
 \sup_{p\in P}\norm{g_{j,p}-g_{j-1,p}}_q
 \leq C\bigl[(1+L_j)2^{2n_{j-1}-aL_j}+2^{-n_{j-1}/2}\bigr],
\end{equation}
for the stated endpoint $q$. The initial filled density is bounded.
Therefore the unmasked stage densities are uniformly Cauchy in $L^q$.
Their weak limits are the actual $\nu_p$, since $D_p$ is continuous
and $\mu_j\rightharpoonup\mu$. On the common bounded output interval,
$L^q$ convergence implies $L^1$ convergence and identifies the density.
No discontinuous mask remains in this final identification.

For $q=\infty$, the unmasked finite densities are jointly continuous
by~\eqref{eq:polar-continuity} with the band indicator removed.
The continuous versions satisfy the same supremum estimates and converge
uniformly. The limit is jointly continuous. In particular the case
$\alpha=3/2$ follows from the explicit summable bound $C2^{-n/2}$
in~\eqref{eq:infinity-endpoint}.
\end{proof}

\subsection{The track lower bound for every horizontal-support pin}
Let $p\in\mathcal B$. At stage $j$ select a supporting horizontal
coding interval of depth $n+L/2$ containing $p_1$. Its width is
$2^{-n-L/2}$ and its exact horizontal mass is
\begin{equation}\label{eq:track-mass}
 m_j=2^{k-bn-aL}.
\end{equation}
This remains valid at coding endpoints by selecting one supporting closed
cylinder. The vertical support has at most $2^{b(n+L)-k}$ occupied
intervals of length $\delta_j=2^{-(n+L)}$. On the horizontal track,
\[
 0\leq|y-p|-(y_2-p_2)
 =\frac{(y_1-p_1)^2}{|y-p|+(y_2-p_2)}
 \leq\frac{2^{-2n-L}}{2v_0}\leq C\delta_j.
\]
Its distance image is therefore contained in a Borel set $U_j$ with
\begin{equation}\label{eq:track-image}
 |U_j|\leq C2^{-c(n+L)-k},\qquad \nu_p(U_j)\geq m_j.
\end{equation}
The argument uses only $p_1\in H$ and the vertical separation, not
membership of $p_2$ in the lower vertical Cantor support.

For every finite $q>1$, H\"older's inequality on this particular set gives
\begin{equation}\label{eq:track-lower}
 \int_{U_j}g_p(r)^q\,dr
 \geq\frac{m_j^q}{|U_j|^{q-1}}
 \geq c_q 2^{[q(c-a)-c]L+[q(c-b)-c]n}.
\end{equation}
The inequality also holds with an infinite left side.
When $1<\alpha<3/2$ and $q>q_*$, the coefficient of $L$ is positive;
superlacunarity makes the right side diverge. Hence $g_p\notin L^q$.
Compact output support then also excludes $L^\infty$.
At $q=q_*$ the residual exponent is instead $-q_*n/2$, consistent with
the attained upper endpoint. This proves the sharp part of
Theorem~\ref{thm:main-super}; off-$H$ behavior is proved in
Section~\ref{sec:pins}.

Consequently the superlacunary regimes on $\mathcal B$ are:
$1<\alpha<4/3$ gives an attained $q_*<2$ and failure of $L^2$;
$\alpha=4/3$ gives $L^2$ and failure of all $L^{2+\varepsilon}$;
$4/3<\alpha<3/2$ gives an attained finite $q_*>2$; and
$\alpha\geq3/2$ gives bounded jointly continuous densities on the entire
pin rectangle. These nonmembership statements never extend to off-$H$ pins.

\section{Fixed ratio and the excluded critical endpoint}\label{sec:fixed}
Assume~\eqref{eq:fixed-schedule}. The fixed-mask proof and interval theorem
already hold when $aK>1$. In this section use the convenient common
hypothesis $aK>2$. Set
\begin{equation}\label{eq:DK}
 D_K=K(c-a)+(c-b)=K(3/2-\alpha)+1-\alpha.
\end{equation}
If $D_K>0$, define $q_K=c(K+1)/D_K$. Notice that
\[
 c(K+1)-D_K=b+aK>0,
\]
so $q_K>1$ as required for the H\"older argument.

\subsection{Upper estimates}
The outer estimate~\eqref{eq:outer} becomes
\begin{equation}\label{eq:fixed-outer}
 C(1+Kn)2^{-(aK-2)n},
\end{equation}
which is summable since $aK>2$ and $n=n_1(K+1)^{j-2}$.
The global inner estimate~\eqref{eq:inner-general} becomes
\begin{align}
 C2^{[-b-aK+c(K+1)(1-1/q)]n}
 &=C2^{[D_K-c(K+1)/q]n}.\label{eq:fixed-inner}
\end{align}
Its constants arise from the same aligned-bin mass bound, true product
structure, and vertical coarea estimate; the old filled law has the
additional favorable factor $2^{-c(1-1/q)L}$.

If $D_K>0$, the exponent in~\eqref{eq:fixed-inner} is negative exactly
when $q<q_K$. If $D_K<0$, it is negative for $q=\infty$. If $D_K=0$,
it is negative for every fixed finite $q$. Combining the inner and outer
bounds proves uniform $L^q$ Cauchy convergence in precisely the upper
ranges asserted in Theorem~\ref{thm:main-fixed}. Weak convergence identifies
the limiting original law. In the $L^\infty$ case the continuous finite
versions converge uniformly, yielding joint continuity.

For a particular finite exponent one can weaken the outer condition.
Every conditional output from a parent square lies in an interval of
length at most $\sqrt2h$, because distance is $1$-Lipschitz. Apply this
support-length gain before summing conditional increments with their
parent masses. Minkowski's inequality gives the outer $L^q$ estimate
\begin{equation}\label{eq:improved-fixed-outer}
 C_q(1+Kn)2^{(2-1/q-aK)n}.
\end{equation}
Thus sufficient conditions for a given finite $q$ are
\begin{equation}\label{eq:weaker-fixed}
 aK>2-1/q,\qquad D_K-c(K+1)/q<0.
\end{equation}
If $D_K>0$, the whole subcritical range follows under the weaker strict
condition $aK>2-1/q_K$. These are sufficient conditions furnished by the
proof, not claims that $aK=1$ or $aK=2$ is an exact geometric threshold.

\subsection{Failure at the endpoint}
For every $p\in\mathcal B$, the track image~\eqref{eq:track-image}
satisfies $|U_j|\to0$. Substituting $L=Kn$ into
\eqref{eq:track-lower} gives
\begin{equation}\label{eq:fixed-lower}
 \int_{U_j}g_p(r)^q\,dr
 \geq c_q 2^{[qD_K-c(K+1)]n}.
\end{equation}
For $q>q_K$ the right side diverges. More importantly, for $q=q_K$
it is bounded below by a fixed positive constant. If $g_p\in L^{q_K}$,
the absolute continuity of the Lebesgue integral of $g_p^{q_K}$ would
force the left side to tend to zero as $|U_j|\to0$, a contradiction.
The $U_j$ need not be disjoint or nested. A constant lower bound for the
\emph{whole} norm alone would not prove endpoint failure; localization to
shrinking-measure sets is essential. This proves the exact open range on
$\mathcal B$.

\subsection{Examples and the endpoint switch}
For $\alpha=13/10$, $a=3/20$ and $b=13/20$. The fixed ratio $K=14$
satisfies $aK=21/10>2$ and gives
\[
 D_K=5/2,\qquad q_K=21/10,\qquad\beta_K=55/31.
\]
All pins in $P$ have $L^2$ original densities, and pins in $\mathcal B$
have exactly $1\leq q<21/10$. For the superlacunary construction with
the same $\alpha$, $q_*=7/4$ is attained, every pin in $\mathcal B$
fails $L^2$, and the planar packing dimension is two. These are two
specific schedules with controlled placements; the comparison does not
show that packing dimension alone determines the phase diagram.

The boundary $D_K=0$ occurs at
\begin{equation}\label{eq:fixed-boundary-alpha}
 \alpha=\frac32-\frac1{2(K+1)}.
\end{equation}
For example $K=9$, $\alpha=29/20$ gives $aK=81/40>2$.
The next section proves noncontinuity on $\mathcal B$ and the logarithmic
norm upper bound there, while keeping the $L^\infty$ question open.

The endpoint contrast has a precise exponent explanation. In the
superlacunary regime the leading critical $L$ exponent cancels while the
inherited horizontal mass leaves the summable $2^{-n/2}$ gain.
At a fixed ratio the shifted critical exponent cancels the entire
$n$-and-$L$ power simultaneously in~\eqref{eq:fixed-lower}, leaving a
positive amount of $L^{q_K}$ mass on the shrinking track images.

\section{Exact exceptional pin sets and the boundary}\label{sec:pins}
\subsection{Off the horizontal support}
\begin{proposition}\label{prop:offH}
Under the superlacunary assumptions, or the fixed-ratio assumptions with
$aK>1$, every $p\in P\setminus\mathcal B$ has a bounded continuous
original density. The density is jointly continuous locally on
$(P\setminus\mathcal B)\times\R$, and
\begin{equation}\label{eq:off-H}
 \norm{g_p}_\infty\leq\frac{C}{\dist(p_1,H)}.
\end{equation}
\end{proposition}
\begin{proof}
Since $H$ is compact, $d=\dist(p_1,H)>0$. Set $\varepsilon=d/2$.
On a neighborhood of $p$ with $\dist(p'_1,H)>\varepsilon$, the mask
$M_\varepsilon(p',\cdot)$ is identically one on the entire source
support. The jointly continuous positive-filter density of
Proposition~\ref{prop:fixed-filter} is therefore the original density
throughout that neighborhood, with bound $2C_*/d$.
\end{proof}
Every fixed off-$H$ pin has a positive gap, so no third membership class
arises from a rate of approximation to $H$. Bounds may diverge as pins
approach $H$, which is a quantitative issue distinct from membership.

\subsection{Dimensions of the full horizontal-support pin set}
\begin{proposition}\label{prop:B-dimensions}
In the superlacunary regime,
\begin{equation}\label{eq:B-super-dim}
 \dimH H=\alpha-1,\quad\dimP H=1,\qquad
 \dimH\mathcal B=\alpha,\quad\dimP\mathcal B=2.
\end{equation}
In the fixed-ratio regime, put
\begin{align}
 s_K&=\frac{b+aK}{1+K/2}
     =\frac{\alpha+(\alpha-1)K}{K+2},\label{eq:sK}\\
 t_K&=\frac{b+aK}{1+aK}
     =\frac{\alpha+(\alpha-1)K}{2+(\alpha-1)K}.\label{eq:tK}
\end{align}
Then
\begin{equation}\label{eq:B-fixed-dim}
 \dimH H=s_K,\quad\dimP H=t_K,\qquad
 \dimH\mathcal B=1+s_K,\quad\dimP\mathcal B=1+t_K.
\end{equation}
\end{proposition}
\begin{proof}
In the superlacunary case,~\eqref{eq:horizontal-frostman} gives the
Hausdorff lower bound for $H$. At depth $n+L/2$, the horizontal free
count is $bn+aL$, whose ratio to depth tends to $2a=\alpha-1$.
The corresponding covers give the upper bound. At depth $n+aL$ the
ratio tends to one, and every fixed cylinder has upper box dimension one.
The coding-cylinder and Baire argument of Section~\ref{sec:construction}
then gives packing dimension one.

In the fixed-ratio case the ratio
$[bn+f_X(l)]/(n+l)$ increases up to $l=aKn$, decreases up to $l=Kn/2$,
and increases thereafter. Its trough is $s_K$ and its peak is $t_K$;
indeed $s_K\leq b\leq t_K$, so no stage endpoint has been omitted.
The finitely many initial depths alter only a constant. The horizontal
law is $s_K$-Frostman, and covers at the middle-stage trough give the
matching Hausdorff upper bound. Every fixed cylinder has upper box
dimension $t_K$, and Baire gives the packing lower bound.

For $\mathcal B=H\times[0,\eta]$, multiply the appropriate horizontal
Frostman law by Lebesgue probability on the interval to obtain the
Hausdorff lower bounds. Multiply the trough covers by interval grids to
obtain the upper bounds. For packing dimension, cylinder products with
nondegenerate subintervals have upper box dimensions larger by exactly one;
the same countable-cover/Baire argument applies. This proves the product
formulas directly rather than relying on an equality for arbitrary
fractal products.
\end{proof}
Useful identities for this family are
\begin{equation}\label{eq:dimension-identities}
 D_K=(K+2)(1/2-s_K),\qquad
 q_K=\frac{1-s_K}{1-2s_K}\quad(D_K>0).
\end{equation}
The superlacunary formula has the same expression in $s=\alpha-1$,
although the endpoint behavior differs.
At $D_K=0$,
\begin{equation}\label{eq:boundary-dimensions}
 s_K=1/2,\quad t_K=\frac{K+1}{K+2},\qquad
 \dimH\mathcal B=3/2,\quad\dimP\mathcal B=2-\frac1{K+2}.
\end{equation}
For $K=9$, $\alpha=29/20$, these dimensions are $3/2$ and $21/11$.
Under the standing requirement $aK>2$, an integer boundary example must
have $K\geq9$.

\subsection{No continuous version at the fixed-ratio boundary}
\begin{proposition}\label{prop:boundary-discontinuity}
If $aK>2$ and $D_K=0$, then every $p\in\mathcal B$ has no continuous
original distance density on $\R$.
\end{proposition}
\begin{proof}
The minimum vertical source coordinate is $r_0=1-\eta$. For
$p_1\in H$, the point $(p_1,r_0)$ lies in the source support, so
\[
 r_{\min}(p)=r_0-p_2>0
\]
is the minimum source distance. Let $\delta_j=2^{-(n+L)}$.
Combine the horizontal track in~\eqref{eq:track-mass} with the lowest
occupied vertical interval of length $\delta_j$, which has
$\zeta_+$-mass $2^{k-b(n+L)}$. The product subset has mass
\begin{equation}\label{eq:boundary-mass}
 2^{2k-2bn-(a+b)L}.
\end{equation}
Its vertical distances are within $\delta_j$ of $r_{\min}(p)$ and its
horizontal quadratic error is at most $C2^{-2n-L}\leq C\delta_j$.
The actual distance image lies in
$[r_{\min}(p),r_{\min}(p)+C_0\delta_j]$. Dividing
\eqref{eq:boundary-mass} by $\delta_j$ gives
\begin{equation}\label{eq:boundary-ratio}
 2^{2k}2^{(1-\alpha)n+(3/2-\alpha)L}=2^{2k}2^{D_Kn}.
\end{equation}
At $D_K=0$ this is a fixed positive constant, so
\begin{equation}\label{eq:boundary-local}
 \nu_p([r_{\min}(p),r_{\min}(p)+C_0\delta_j])\geq c_0\delta_j.
\end{equation}
A continuous density for $\nu_p$ must vanish to the left of its support
and hence at $r_{\min}(p)$. Its supremum on the shrinking displayed
interval would tend to zero, giving mass $o(\delta_j)$, contrary to
\eqref{eq:boundary-local}.
\end{proof}
This proof does not exclude an essentially bounded density with a jump at
the support endpoint. It settles continuity, not $L^\infty$ membership.

\subsection{The exact classification}
Define
\begin{align*}
 \Bad_q&=\{p\in P:\nu_p\text{ has no }L^q\text{ density}\},\\
 \Bad_{C^0}&=\{p\in P:\nu_p\text{ has no continuous density on }\R\}.
\end{align*}
All these laws are absolutely continuous. Compact support implies that
every continuous-density version is bounded. The upper bounds,
\eqref{eq:track-lower},~\eqref{eq:fixed-lower},
Proposition~\ref{prop:offH}, and Proposition~\ref{prop:boundary-discontinuity}
therefore give the following exact statements.
\begin{enumerate}[leftmargin=*]
\item Superlacunary, $1<\alpha<3/2$:
\[
 \Bad_q=\varnothing\ (1\leq q\leq q_*),\qquad
 \Bad_q=\mathcal B\ (q_*<q\leq\infty),\qquad
 \Bad_{C^0}=\mathcal B.
\]
\item Superlacunary, $3/2\leq\alpha<2$: all $\Bad_q$,
$1\leq q\leq\infty$, and $\Bad_{C^0}$ are empty.
\item Fixed ratio, $aK>2$ and $D_K>0$:
\[
 \Bad_q=\varnothing\ (1\leq q<q_K),\qquad
 \Bad_q=\mathcal B\ (q_K\leq q\leq\infty),\qquad
 \Bad_{C^0}=\mathcal B.
\]
\item Fixed ratio, $aK>2$ and $D_K<0$: all these exceptional sets
are empty.
\item Fixed ratio, $aK>2$ and $D_K=0$:
\[
 \Bad_q=\varnothing\ (1\leq q<\infty),\qquad
 \Bad_{C^0}=\mathcal B,\qquad \Bad_\infty\subseteq\mathcal B.
\]
The exact set $\Bad_\infty$ in the last line is not determined here.
\end{enumerate}
The nonempty exact exceptional sets have the dimensions in
Proposition~\ref{prop:B-dimensions}. They are exceptions to stronger
\emph{density regularity}, not pins with zero-measure or interval-free
distance sets. Every pin in $P$ still has the interval conclusion.

\subsection{An upper bound on all boundary norms}
\begin{theorem}\label{thm:log-upper}
Assume $aK>2$ and $D_K=0$. Then for a fixed $C<\infty$,
\begin{equation}\label{eq:log-upper}
 \sup_{p\in P}\norm{g_p}_q\leq C\log(e+q)
 \qquad(1\leq q<\infty).
\end{equation}
There are $c_1,c_2,T_0>0$ such that
\begin{equation}\label{eq:double-tail}
 \sup_{p\in P}|\{r:g_p(r)>T\}|
 \leq \exp[-c_1\exp(c_2T)]\qquad(T\geq T_0).
\end{equation}
For some $c_3>0$,
\begin{equation}\label{eq:double-integrability}
 \sup_{p\in P}\int_I\exp(\exp(c_3g_p(r)))\,dr<\infty.
\end{equation}
Only upper bounds are asserted; no matching lower growth in $q$ or tail
asymptotic is proved.
\end{theorem}
\begin{proof}
The initial density and the summable outer increments
\eqref{eq:fixed-outer} have $L^q$ bounds independent of $q\geq1$, using
the fixed finite output interval. The constant in~\eqref{eq:MH} can also
be chosen uniformly in $q$: the coarea factor is
$(D_0/v_0)^{1-1/q}$ and fixed prefix factors are bounded by their
$q=\infty$ values. At $D_K=0$, the inner increment is at most
$C2^{-c(K+1)n/q}$. Thus
\begin{equation}\label{eq:lacunary-moment-series}
 \norm{g_p}_q\leq C+C\sum_{j\geq2}
 2^{-c(K+1)n_1(K+1)^{j-2}/q}.
\end{equation}
There are $O(1+\log q)$ terms before the exponent is of order one;
after that point the remaining sum is bounded by a fixed convergent
series of the form $\sum_{l\geq0}\exp[-A(K+1)^l]$.
This proves~\eqref{eq:log-upper}.

For large $T$ choose the real moment order $q=\exp[T/(4C)]$,
increasing $C$ and the lower threshold if necessary. Then
$C\log(e+q)\leq T/2$ and Markov gives
\[
 |\{g_p>T\}|\leq(\norm{g_p}_q/T)^q\leq2^{-q},
\]
which proves~\eqref{eq:double-tail}. Choose $0<c_3<c_2$ sufficiently
small. Layer-cake integration against the derivative of
$\exp(\exp(c_3T))$ gives~\eqref{eq:double-integrability}, since the
negative term $-c_1\exp(c_2T)$ dominates $\exp(c_3T)$ and the remaining
linear terms. The integral is over $I$, not all of $\R$, so its nonzero
constant baseline causes no divergence.
\end{proof}
The double-exponential upper tail neither proves nor disproves essential
boundedness. This completes all assertions of Theorem~\ref{thm:main-fixed}
and the full pin classification, with the stated boundary question intact.

\section{Positive masks and nonuniform replacement stability}\label{sec:stability}
The fixed restriction controls arbitrary further positive source masks
without requiring cancellation, a product structure for those masks, or
an independent copy of the pin.

\subsection{Domination and the retained mass cost}
Fix $\varepsilon>0$, and let $0\leq F(p,y)\leq M_\varepsilon(p,y)$
be jointly Borel. Set
\[
 \sigma_p=(D_p)_*(F(p,\cdot)\mu),\qquad
 m(p)=\int F(p,y)\,d\mu(y).
\]
Positive domination gives a density $f_p$ satisfying
\begin{equation}\label{eq:domination}
 0\leq f_p\leq g_\varepsilon(p,\cdot)\leq C_*/\varepsilon.
\end{equation}
Consequently
\begin{equation}\label{eq:mask-L2}
 \norm{f_p}_2^2\leq(C_*/\varepsilon)m(p).
\end{equation}
For $m(p)>0$, normalization yields
\begin{equation}\label{eq:normalized}
 \norm{f_p/m(p)}_\infty\leq\frac{C_*}{\varepsilon m(p)},
 \qquad
 \norm{f_p/m(p)}_2^2\leq\frac{C_*}{\varepsilon m(p)}.
\end{equation}
The squared $L^2$ norm is divided by $m(p)^2$; the unnormalized bound
already contains one factor $m(p)$. Uniform normalized bounds therefore
require a genuine lower bound on retained mass. One may integrate these
fiber inequalities against any probability on $P$, using that same pin
in both distance factors. Jointly measurable versions follow from the
finite coarea formulas and convergence in the separable space $L^1(I)$.

Whenever the original density itself is uniformly bounded, as in the
superlacunary range $\alpha\geq3/2$ and the fixed-ratio range $D_K<0$
under the stated growth hypothesis, the same argument applies to every
$0\leq F\leq1$ without a preliminary band, using the raw supremum bound.

\subsection{Previously imposed decreasing restrictions}
Suppose $0\leq\Gamma_N\leq1$ is a decreasing sequence and
\[
 \int\Gamma_\infty\,d(\lambda\times\mu)\geq c_*>0.
\]
Choose one $\varepsilon$ with $K_X\varepsilon^\beta<c_*/2$ and put
$F_N=M_\varepsilon\Gamma_N$. This is still a decreasing family,
loses at most $K_X\varepsilon^\beta$ additional pair mass, and has
limiting pair mass at least $c_*/2$. Every fiber has the density bound
\eqref{eq:domination}; the limiting family has nonzero fibers on a
positive-$\lambda$ set of pins. All conditions encoded by $\Gamma_N$
remain imposed.

This multiplication need not preserve a previously asserted relative
retained fraction in every spatial cell. Such local fractions must be
checked separately or obtained by further pruning with their actual mass
budget. Positive domination survives any further pruning but does not
supply its missing mass.

\subsection{Endpoint perturbations and whole low frequency bounds}
For a positive measure $\tau$ on $\R$, define
\[
 \operatorname{Coll}_\delta(\tau)
 =\iint\ind_{\{|s-t|\leq\delta\}}\,d\tau(s)\,d\tau(t).
\]
Move a retained source endpoint by at most $A_0\delta$, possibly through
a pin-dependent probability kernel. Pull the retained coupling back to its
original marginal $F(p,y)\mu$ with $F\leq M_\varepsilon$ and mass
$m(p)$. If two perturbed distances differ by at most $\delta$, their
original distances differ by at most $(1+2A_0)\delta$. Since the original
positive density is bounded by $C_*/\varepsilon$, the perturbed law
$\widetilde\sigma_p$ satisfies
\begin{equation}\label{eq:perturbed-collision}
 \operatorname{Coll}_\delta(\widetilde\sigma_p)
 \leq2(1+2A_0)\delta\,(C_*/\varepsilon)m(p).
\end{equation}
After division by $m(p)^2$, the probability version is bounded by
$2(1+2A_0)\delta C_*/[\varepsilon m(p)]$.
A shared pin displacement of magnitude at most $A_p\delta$ changes the
factor $1+2A_0$ to $1+2A_0+2A_p$. The same physical pin is used before
and after each two-distance comparison.

For completeness, with Fourier convention
$\widehat\tau(\xi)=\int e^{-2\pi i\xi t}\,d\tau(t)$, positivity gives
\begin{equation}\label{eq:lowpass}
 \int_{-R}^R|\widehat\tau(\xi)|^2\,d\xi
 \leq C R\operatorname{Coll}_{1/R}(\tau).
\end{equation}
Indeed take $k_\delta(t)=(1-|t|/\delta)_+$ with $\delta=(2R)^{-1}$.
Its Fourier transform is
$\delta[\sin(\pi\delta\xi)/(\pi\delta\xi)]^2$, nonnegative and
bounded below by a positive multiple of $\delta$ on $[-R,R]$.
Fourier inversion gives
\[
 \int\widehat k_\delta(\xi)|\widehat\tau(\xi)|^2\,d\xi
 =\iint k_\delta(s-t)\,d\tau(s)\,d\tau(t)
 \leq\operatorname{Coll}_{1/R}(\tau),
\]
which proves~\eqref{eq:lowpass}. Finite-measure Fourier inversion is
legitimate since $\widehat k_\delta$ is integrable.
Thus~\eqref{eq:perturbed-collision} supplies whole-low-pass control.
It does not imply a decaying high-frequency shell estimate for arbitrary
varying masks or endpoint perturbations.

\subsection{A weighted nonuniform sufficient condition}
The positive-filter argument has a broader deterministic form that does
not require a globally product limiting source. Its conclusion is weaker
than the exact endpoint phase diagrams.

\begin{theorem}[Weighted rectangular replacements]\label{thm:nonuniform}
Let $\rho_j$ be probabilities on finite unions of filled dyadic squares
in the separated source rectangle, converging weakly to $\rho$.
At step $j$, replace each parent square $Q$, of side $h_Q$ and mass $m_Q$,
by the affine product template of Section~\ref{sec:quadrature}, preserving
its mass. Allow the parameters $A_Q,B_Q,w_Q,v_Q$ and masses to depend
on $Q$. Suppose the new filled leaves refine their parents, their maximum
diameter tends to zero, and
\begin{equation}\label{eq:nonuniform-E}
 \mathcal E=\sum_j\sum_{Q\text{ at step }j}\frac{m_Q}{h_Q}
 \left(A_Q^{-1}+(B_Qw_Q)^{-1}\right)<\infty.
\end{equation}
Assume the first-coordinate marginal of $\rho$ has no atoms. Then for
each fixed $\varepsilon>0$ and every pin in the separated pin rectangle,
$(D_p)_*(M_\varepsilon\rho)$ has a density bounded by
$C(1+\varepsilon^{-1})$. Every original $(D_p)_*\rho$ is absolutely
continuous. If the horizontal marginal satisfies a $\gamma$-Frostman
interval bound, the raw densities are uniformly weak-$L^{1+\gamma}$.
\end{theorem}
\begin{proof}
Apply~\eqref{eq:fixed-band} with each parent's parameters, multiply by
$m_Q$, and sum. Condition~\eqref{eq:nonuniform-E} makes the density
increments absolutely summable in $L^\infty$. The initial finite mixture
of filled squares has a bounded distance density by vertical coarea.
Nonatomicity makes both hard-cut boundary lines $\rho$-null, so the same
weak-limit identification as in Proposition~\ref{prop:fixed-filter} applies.
Increasing the positive restrictions as $\varepsilon\downarrow0$ gives
raw absolute continuity. Finally the lost mass is $O(\varepsilon^\gamma)$;
the Markov decomposition in Corollary~\ref{cor:weak} yields the
weak-$L^{1+\gamma}$ estimate.
\end{proof}
For uniform stages~\eqref{eq:nonuniform-E} reduces to
\[
 \sum_j 2^{n_{j-1}}
       \bigl[A_j^{-1}+(B_jw_j)^{-1}\bigr]<\infty.
\]
The criterion records actual spatial resolution and accumulated weighted
error. It covers adaptive parent-dependent templates with a common
coordinate orientation, not arbitrary rotations. It does not by itself
imply the exact critical exponents: those also use the globally aligned
bin masses and product marginals in~\eqref{eq:global-M}--\eqref{eq:global-H}.

\section{Scope and remaining questions}\label{sec:scope}
All conclusions concern the deterministic, rational-parameter digit models
specified here and their original natural measures. The actual rectangular
placement, the resolving relation $Bw=A$, and the schedule are substantive
hypotheses. Neither the Frostman exponent, the packing dimension, nor a
branch-count profile alone supplies them. The proof uses elementary bounded
variation, coarea in a strictly monotone coordinate, weak and dominated
convergence, interpolation, and H\"older and Markov inequalities. No general
projection theorem, random-martingale theorem, or unrestricted recursive
common-pin estimate is an input.

Fixing a sufficiently small $\varepsilon$ gives an arbitrarily high
retained fraction with a fixed good-law bound. Sending
$\varepsilon\downarrow0$ restores the original law, but the available
bound $C_*/\varepsilon$ diverges. In the superlacunary range
$1<\alpha<4/3$, uniform $L^2$ restoration at a pin in $\mathcal B$
would contradict the proved non-$L^2$ original law: an $L^2$-bounded
sequence has a weak $L^2$ limit representing any weakly convergent measure
limit. The auxiliary changing bands used to estimate signed increments
were never assigned such a uniform $L^2$ bound.

The following matters remain outside the proved conclusions.
\begin{enumerate}[leftmargin=*]
\item At $D_K=0$ and $p_1\in H$, essential boundedness and the exact
set $\Bad_\infty$ remain unresolved. The noncontinuity argument allows a
bounded jump, and the norm upper bound does not force unboundedness.
\item No matching lower bound for the $O(\log q)$ boundary norms or
sharp double-exponential tail asymptotic is proved. At the finite
fixed-ratio endpoint, weak-$L^{q_K}$ or a sharper Lorentz classification
is not determined.
\item The inequalities $aK>1$ and $aK>2$ are sufficient conditions
for the displayed estimates. Equality cases, smaller ratios, and an exact
necessary growth threshold are not settled.
\item No loss-free extension to arbitrary irrational $\alpha$, arbitrary
rotated templates, or exact nonproduct endpoints is asserted. Arbitrary
source masks inserted \emph{before} a replacement comparison need not
preserve its geometry; positive domination applies to masks of the actual
fixed good law after its construction.
\item Local mass retention for unrelated pre-existing pruning procedures
must be checked with their own budgets. The estimates do not establish
an unrestricted common-pin recursion or arbitrary-mask high-frequency
shell decay.
\end{enumerate}
The classical finite-scale insight is attributed to~\cite{GIOW}; the
model-specific conclusions here are not advertised as a new general
Falconer threshold. A broad novelty assessment, professional peer review,
and formal proof verification are not supplied by this manuscript.

\appendix
\section{A direct subcritical increment argument}\label{app:subcritical}
The endpoint proof in Section~\ref{sec:superlacunary} uses a harmonic gain
and global product masses. We retain here the alternative strictly
subcritical argument, which uses a wider auxiliary band and estimates
inside each parent square. Assume $1<\alpha<3/2$ and fix $1<q<q_*$.
Set
\begin{equation}\label{eq:theta}
 \sigma=1-1/q,\qquad
 \theta=\frac{1+c\sigma/a}{2},\qquad
 \gamma_q=\frac{a-c\sigma}{2}>0.
\end{equation}
Then $0<\theta<1$. In a parent square of side $h=2^{-n}$, take the
physical band width $\varepsilon_j=hd$ where $d=A^{-\theta}$.
This is again a decomposition of a signed increment, not a good-law
restoration with a uniform $L^2$ bound.

By Lemma~\ref{lem:fixed-band} and $Bw=A$, the outer conditional
$L^\infty$ error is at most
\[
 Ch^{-1}\bigl[A^{-1}+(hdA)^{-1}\bigr]
 \leq Ch^{-2}A^{-(1-\theta)}.
\]
The output lies in an interval of length $\sqrt2h$, so its $L^q$ error
is at most
\begin{equation}\label{eq:subcritical-outer}
 C_qh^{-(2-1/q)}A^{-(1-\theta)}.
\end{equation}
Inside the band, the replacement has mass at most $2d+2/A\leq4d$,
since $d\geq1/A$ and the horizontal law assigns mass $1/A$ to each
bin. The filled parent has band mass at most $2d$. The replacement's
vertical relative-coordinate density is at most $(Bv)^{-1}$; physical
vertical coarea gives an inner density supremum at most
$Ch^{-1}d(Bv)^{-1}$. Its mass is at most $4d$, so interpolation yields
\begin{equation}\label{eq:subcritical-inner}
 \norm{\text{inner conditional density}}_q
 \leq C_q h^{-\sigma}d(Bv)^{-\sigma}.
\end{equation}
The parent obeys the same bound with $Bv$ replaced by one and is therefore
no larger. These finite-stage density estimates do not presume a
Lebesgue vertical coordinate in the limiting singular law.

Using $Bv=2^{-cL}$, combine~\eqref{eq:subcritical-outer} and
\eqref{eq:subcritical-inner}. The conditional unmasked increment is at most
\begin{align}
 C_q\bigl[h^{-(2-1/q)}2^{-a(1-\theta)L}
       +h^{-\sigma}2^{-(a\theta-c\sigma)L}\bigr]
 &\leq C_q2^{(2-1/q)n-\gamma_qL},\label{eq:subcritical-increment}
\end{align}
since both $L$ exponents equal $\gamma_q$ by~\eqref{eq:theta}.
Weighting by parent masses and applying Minkowski leaves this bound
unchanged. Superlacunarity makes the series summable; for example,
eventually $\gamma_qL_j\geq(3-1/q)n_{j-1}$, leaving $C_q2^{-n_{j-1}}$.
The initial density is bounded and weak convergence identifies the $L^q$
limit with the actual original law. This recovers the strictly subcritical
upper range directly. The separate harmonic and global-band argument is
what retains the finite endpoint and proves the bounded endpoint at
$\alpha=3/2$.

\begin{thebibliography}{9}
\bibitem{GIOW}
L.~Guth, A.~Iosevich, Y.~Ou, and H.~Wang,
\emph{On Falconer's distance set problem in the plane},
arXiv:1808.09346v1 (2018), especially Sections~1.2 and~6.
\url{https://arxiv.org/abs/1808.09346v1}.
\bibitem{PS}
Y.~Peres and W.~Schlag,
\emph{Smoothness of projections, Bernoulli convolutions, and the dimension
of exceptions}, Duke Mathematical Journal \textbf{102} (2000), 193--251.
\url{https://doi.org/10.1215/S0012-7094-00-10222-0}.
\end{thebibliography}
\end{document}
